% Bagian: computability
% Bab: computability-theory
% Subbagian: ce-sets

\documentclass[../../../include/open-logic-section]{subfiles}

\begin{document}

\olfileid[id]{cmp}{thy}{ces}
\olsection{Himpunan Terenumerasi Secara Komputabel}

\begin{defn}
Suatu himpunan \emph{terenumerasi secara komputabel} jika himpunan itu kosong
atau merupakan daerah hasil suatu fungsi dapat dihitung.
\end{defn}

\begin{history}
Himpunan terenumerasi secara komputabel juga disebut \emph{terenumerasi secara
  rekursif}. Istilah terakhir merupakan istilah asli, dan saat ini keduanya
lazim digunakan, begitu pula singkatan ``c.e.'' dan ``r.e.''
\end{history}

\begin{explain}
Anda sebaiknya memikirkan arti definisi itu dan alasan istilah tersebut tepat.
Gagasannya adalah bahwa jika $S$ merupakan daerah hasil fungsi dapat
dihitung~$f$, maka
\[
S = \{ f(0), f(1), f(2), \dots \},
\]
sehingga $f$ dapat dipandang ``mengenumerasi'' unsur-unsur~$S$. Perhatikan
bahwa berdasarkan definisi, $f$ tidak harus berupa fungsi naik; yakni,
enumerasinya tidak harus berurutan menaik. Bahkan, $f$ tidak harus injektif;
artinya, pengulangan dalam enumerasi $f(0)$, $f(1)$, $f(2)$, \dots{} atas~$S$
diperbolehkan. Sebagai contoh, fungsi konstan $f(x) = 0$ mengenumerasi
himpunan $\{ 0 \}$.
\end{explain}

Setiap himpunan dapat dihitung terenumerasi secara komputabel. Untuk melihat
hal ini, andaikan $S$ dapat dihitung. Jika $S$ kosong, berdasarkan definisi
himpunan itu terenumerasi secara komputabel. Jika tidak, misalkan $a$ sembarang
unsur $S$. Definisikan $f$ dengan
\[
f(x) =
\begin{cases}
x & \text{jika $\Char{S}(x) = 1$} \\
a & \text{selain itu.}
\end{cases}
\]
Maka $f$ merupakan fungsi dapat dihitung dan $S$ merupakan daerah hasil~$f$.

\end{document}

